% Copyright 2018 Daniel Nikpayuk
\documentclass[twoside]{article}
\usepackage[letterpaper,left=1cm,right=1cm,top=2cm,bottom=2cm]{geometry}
\usepackage{amsfonts}
\usepackage{graphicx}
\usepackage{hyperref}
\usepackage{amsmath}

\newtheorem{theorem}{Theorem}
\newtheorem{lemma}[theorem]{Lemma}
\newtheorem{proposition}[theorem]{Proposition}
\newtheorem{corollary}[theorem]{Corollary}

\newenvironment{proof}[1][Proof]{\begin{trivlist}
\item[\hskip \labelsep {\bfseries #1}]}{\end{trivlist}}
\newenvironment{definition}[1][Definition]{\begin{trivlist}
\item[\hskip \labelsep {\bfseries #1:}]}{\end{trivlist}}
\newenvironment{example}[1][Example]{\begin{trivlist}
\item[\hskip \labelsep {\bfseries #1}]}{\end{trivlist}}
\newenvironment{remark}[1][Remark]{\begin{trivlist}
\item[\hskip \labelsep {\bfseries #1}]}{\end{trivlist}}

\newcommand{\qed}{\nobreak \ifvmode \relax \else
      \ifdim\lastskip<1.5em \hskip-\lastskip
      \hskip1.5em plus0em minus0.5em \fi \nobreak
      \vrule height0.75em width0.5em depth0.25em\fi}

\title{Interface Theory}
\author{Daniel Nikpayuk}
\date{March 2, 2018}
\pagestyle{empty}
\begin{document}
\maketitle
\thispagestyle{empty}

\begin{figure}[h]
\centering
\includegraphics[width=1in]{../../../cc-by-nc.png}\\[0.1in]
\tiny This article is licensed under \\
\href{http://creativecommons.org/licenses/by-nc/4.0/}
{Creative Commons Attribution-NonCommercial 4.0 International.}\\[0.3in]
\end{figure}

The general formula for cleanly designing any interface:

\begin{enumerate}
\item Constructively define the structure being modeled.
\item Define every possible substructure.
\item Find all possible algebras (spaces) of substructures.
\item Compress each substructure space, this is an interface.
\item Determine how each interface relates to every other interface.
\item Use chains of interfaces to stratify and mitigate levels of interface privilege,
      this allows for the trade-off between optimization and safety.
\end{enumerate}

This is the direct approach, although in practice if your initial structure is infinite in size, there will be infinite combinatorial
possibilities so it becomes impractical. Even in finite cases, the combinatorial possibilities generally increase rapidly,
so heuristic approaches will still need to be taken. Very rarely can you actually brute force the algorithm,
but just knowing what it is actually helps quite a lot.

Following this, one of the first orders of business is to find a basis if possible. As our entire space is constructive,
what are the constructive grammatical elements that allow us to build any given interface? Finding such bases reveals
the bare minimum we'd need to equip a library with to build a given interface. In practice, a common paradigm or template
strategy is to look at the lifecycle ecology of the constructs within the given interface.

Another first order of business is to prove interface chains are \emph{transitive}, meaning for every subinterface
(indirect interface), there exists an isomorphic direct interface. From a user's perspective, this means you can interface
indirectly, then subinterface, then subinterface, and so on, until you're working with your intended interface resolution,
or you can work directly at that resolution level, or any level in-between. From a compiler's perspective, this means
you can inline (translate) any indirect grammar into direct grammar for optimization purposes.

\end{document}
